\section{Formulario Macroeconómico}

\subsection*{I. Tasas de Crecimiento y Contabilidad Nacional}

\begin{itemize}
	\item \textbf{Tasa de Variación Interanual (TVI):} Mide el crecimiento de un año a otro.
\end{itemize}

\[
TVI = \frac{V_t - V_{t-1}}{V_{t-1}} \cdot 100
\]

\begin{itemize}
	\item \textbf{Tasa de Variación Media Anual Acumulativa (TVMAA):} Calcula el ritmo de crecimiento promedio durante un periodo de $t$ años.
\end{itemize}

\[
TVMAA = \left( \left( \frac{V_t}{V_0} \right)^{\frac{1}{t}} - 1 \right) \cdot 100
\]

\begin{itemize}
	\item \textbf{Deflactor Implícito del PIB:} Mide la evolución de los precios de todos los bienes y servicios de la economía.
\end{itemize}

\[
Deflactor = \frac{PIB_{nominal}}{PIB_{real}} \cdot 100
\]

\begin{itemize}
	\item \textbf{Cálculo del PIB real (a precios constantes):} Elimina el efecto de la inflación.
\end{itemize}

\[
PIB_{real} = \frac{PIB_{nominal}}{Deflactor} \cdot 100
\]

\begin{itemize}
	\item \textbf{PIB desde la Demanda:}
\end{itemize}

\[
PIB = C + FBC + X - M
\]

\textit{(Donde $C$ = Consumo final, $FBC$ = Formación Bruta de Capital, $X$ = Exportaciones, $M$ = Importaciones)}

\begin{itemize}
	\item \textbf{PIB desde la Oferta:}
\end{itemize}

\[
PIB = \sum VAB_{sectores} + Impuestos\,netos\,sobre\,productos
\]

\[
VAB = Producci\'on - Consumo\,intermedio
\]

\begin{itemize}
	\item \textbf{Valor Añadido Bruto a coste de los factores (VABcf):} Es la suma de lo que perciben los trabajadores y los empresarios/propietarios por su actividad.
\end{itemize}

\[
VAB_{cf} = \text{Remuneraci\'on de Asalariados (RA)} + \text{Excedente Bruto de Explotaci\'on/Rentas Mixtas (EBE)}
\]

\begin{itemize}
	\item \textbf{Valor Añadido Bruto a precios b\'asicos (VABpb):} Se obtiene sumando al coste de los factores los impuestos que gravan la producci\'on (restando las subvenciones recibidas).
\end{itemize}

\[
VAB_{pb} = VAB_{cf} + \text{Otros Impuestos Netos sobre la Producci\'on (OINP)}
\]

\begin{itemize}
	\item \textbf{Otros Impuestos Netos sobre la Producci\'on (OINP):} Diferencia entre impuestos que gravan la producci\'on y las subvenciones recibidas.
\end{itemize}

\[
OINP = Impuestos - Subvenciones
\]

\subsection*{II. Mercado de Trabajo (Tema 7 y Cuaderno 4)}

\begin{itemize}
	\item \textbf{Tasa de Actividad:}
\end{itemize}

\[
Tasa\,de\,Actividad = \frac{Poblaci\'on\,Activa}{Poblaci\'on\,\geq\,16\,a\~nos} \cdot 100
\]

\begin{itemize}
	\item \textbf{Tasa de Empleo (u Ocupación):}
\end{itemize}

\[
Tasa\,de\,Empleo = \frac{Poblaci\'on\,Ocupada}{Poblaci\'on\,\geq\,16\,a\~nos} \cdot 100
\]

\begin{itemize}
	\item \textbf{Tasa de Paro (o Desempleo):}
\end{itemize}

\[
Tasa\,de\,Paro = \frac{Poblaci\'on\,Desempleada}{Poblaci\'on\,Activa} \cdot 100
\]

\begin{itemize}
	\item \textbf{Tasa de Asalarización:}
\end{itemize}

\[
Tasa\,de\,Asalarizaci\'on = \frac{N^{\circ}\,Asalariados}{Poblaci\'on\,Ocupada} \cdot 100
\]

\begin{itemize}
	\item \textbf{Tasa de Temporalidad:}
\end{itemize}

\[
Tasa\,de\,Temporalidad = \frac{N^{\circ}\,Asalariados\,con\,contrato\,temporal}{N^{\circ}\,Total\,de\,Asalariados} \cdot 100
\]

\begin{itemize}
	\item \textbf{Tasa de Trabajo a Tiempo Parcial:}
\end{itemize}

\[
Tasa\,a\,tiempo\,parcial = \frac{N^{\circ}\,Ocupados\,a\,tiempo\,parcial}{Poblaci\'on\,Ocupada} \cdot 100
\]

\subsection*{III. Productividad, Competitividad y Crecimiento (Temas 2, 3, 4 y Cuadernos 1, 2)}

\begin{itemize}
	\item \textbf{Descomposición del PIB per c\'apita (Forma detallada con Ocupados):}
\end{itemize}

\[
\frac{\text{PIB}_{real}}{\text{Poblaci\'on}} = \frac{\text{PIB}_{real}}{\text{Ocupados}} \times \frac{\text{Ocupados}}{\text{Poblaci\'on Activa}} \times \frac{\text{Poblaci\'on Activa}}{\text{Poblaci\'on}}
\]

\begin{itemize}
	\item \textbf{Descomposición del PIB per c\'apita (Forma simplificada):}
\end{itemize}

\[
\frac{\text{PIB}}{\text{Poblaci\'on}} = \frac{\text{PIB}}{\text{Empleo}} \times \frac{\text{Empleo}}{\text{Poblaci\'on}}
\]

\begin{itemize}
	\item \textbf{Descomposición de la Tasa Global de Empleo:}
\end{itemize}

\[
\frac{Empleo}{Poblaci\'on} = \frac{Empleo}{PET} \cdot \frac{PET}{Poblaci\'on}
\]

\textit{(Donde $PET$ es la Población en Edad de Trabajar)}

\begin{itemize}
	\item \textbf{Productividad Aparente del Trabajo (PAT):}
\end{itemize}

\[
PAT = \frac{PIB_{real}}{Empleo_{total}}
\]

\textit{(También aplicable sectorialmente como $VAB / Empleo$)}

\begin{itemize}
	\item \textbf{Descomposición de la Productividad Agraria (Tema 3):}
\end{itemize}

\[
\frac{Q}{L} = \frac{Q}{T} \cdot \frac{T}{L}
\]

\textit{(Productividad del trabajo = Productividad de la tierra $\times$ Ratio de estructuras o superficie por ocupado)}

\begin{itemize}
	\item \textbf{Costes Laborales Unitarios Nominales ($CLU_{nominal}$):}
\end{itemize}

\[
CLU_{nominal} = \frac{Remuneraci\'on\,por\,Asalariado}{PAT} = \frac{RA / Asalariados}{PIB_{nominal} / Empleo}
\]

\begin{itemize}
	\item \textbf{Costes Laborales Unitarios Reales ($CLU_{real}$):}
\end{itemize}

\[
CLU_{real} = \frac{CLU_{nominal}}{Deflactor\,del\,PIB}
\]

\begin{itemize}
	\item \textbf{Función de Producción Cobb-Douglas (Productividad):}
\end{itemize}

\[
y = A \cdot k^\alpha \cdot h^\beta
\]

\textit{(Donde $y$ = productividad, $A$ = progreso técnico, $k$ = capital físico por trabajador, $h$ = capital humano)}

\begin{itemize}
	\item \textbf{Residuo de Solow (Cálculo del Progreso Técnico o PTF):} Expresado en tasas de variación.
\end{itemize}

\[
\Delta A = \Delta y - \alpha\Delta k - \beta\Delta h
\]

\subsection*{IV. Sector Exterior y Balanza de Pagos (Tema 6 y Cuaderno 3)}

\begin{itemize}
	\item \textbf{Propensión a Exportar:}
\end{itemize}

\[
P_x = \frac{X}{PIB} \cdot 100
\]

\begin{itemize}
	\item \textbf{Propensión a Importar:}
\end{itemize}

\[
P_m = \frac{M}{PIB} \cdot 100
\]

\begin{itemize}
	\item \textbf{Grado de Apertura Exterior:}
\end{itemize}

\[
Apertura = \frac{X + M}{PIB} \cdot 100
\]

\begin{itemize}
	\item \textbf{Consumo Aparente (CA):} Representa la demanda interna total de los bienes o servicios de un sector. Es decir, lo que gasta el pa\'is en esos productos, independientemente de d\'onde se hayan fabricado.
\end{itemize}

\[
CA = Producci\'on + Importaciones - Exportaciones
\]

\begin{itemize}
	\item \textbf{Grado de Penetraci\'on de las Importaciones (GPI):} Mide qu\'e porcentaje del consumo interno de un pa\'is se satisface con productos comprados al extranjero.
\end{itemize}

\[
GPI = \left( \frac{Importaciones}{CA} \right) \times 100
\]

\begin{itemize}
	\item \textbf{Tasa de Cobertura Comercial:} Porcentaje de las importaciones que se pueden pagar con las exportaciones.
\end{itemize}

\[
Cobertura = \frac{X}{M} \cdot 100
\]

\begin{itemize}
	\item \textbf{Exportaciones Netas (Saldo Comercial Absoluto):}
\end{itemize}

\[
X_n = X - M
\]

\begin{itemize}
	\item \textbf{Saldo Comercial Relativo:}
\end{itemize}

\[
SCR = \frac{X - M}{X + M} \cdot 100
\]

\begin{itemize}
	\item \textbf{Capacidad (+) / Necesidad (-) de Financiación (Visión Real):}
\end{itemize}

\[
Saldo = Saldo\,Cuenta\,Corriente + Saldo\,Cuenta\,Capital
\]

\begin{itemize}
	\item \textbf{Capacidad (+) / Necesidad (-) de Financiación (Visión Financiera):}
\end{itemize}

\[
Saldo = Variaci\'on\,Neta\,de\,Activos\,(VNA) - Variaci\'on\,Neta\,de\,Pasivos\,(VNP)
\]

\subsubsection*{Fórmulas y Resoluciones Paso a Paso de la Balanza de Pagos}

Al igual que en los ejercicios anteriores, en macroeconomía se utilizan identidades contables fundamentales para resolver los huecos de la balanza de pagos. Se calcula ``de abajo hacia arriba'' aprovechando los datos totales.

\begin{itemize}
	\item \textbf{1. Saldo Corriente y Capital (Capacidad o Necesidad de Financiación)}
\end{itemize}

F\'ormula: $\text{Saldo Cuenta Financiera} - \text{Errores y omisiones}$

C\'alculo: $-101.869,0 - 480,1 = -102.349,1$

Nota: Como el resultado es negativo, indica que la econom\'ia tiene una necesidad de financiaci\'on frente al resto del mundo.

\begin{itemize}
	\item \textbf{2. Saldo de la Cuenta Corriente}
\end{itemize}

F\'ormula: $\text{Saldo Corriente y capital} - \text{Saldo de Capital}$

C\'alculo: $-102.349,1 - 4.812,0 = -107.161,1$

\begin{itemize}
	\item \textbf{3. Saldo de Servicios}
\end{itemize}

F\'ormula: $\text{Saldo Corriente} - (\text{Saldo Comercial} + \text{Saldo Rentas Primarias} + \text{Saldo Rentas Secundarias})$

C\'alculo: $-107.161,1 - (-88.459,2 - 32.912,8 - 7.199,9)$

Desarrollo: $-107.161,1 - (-128.571,9) = 21.410,8$

\begin{itemize}
	\item \textbf{4. Pagos por Servicios}
\end{itemize}

F\'ormula: $\text{Ingresos de Servicios} - \text{Saldo de Servicios}$

C\'alculo: $93.404,4 - 21.410,8 = 71.993,6$

\subsection*{V. Sector Público (Tema 9 y Cuaderno 4)}

\begin{itemize}
	\item \textbf{Peso del Gasto Público:}
\end{itemize}

\[
Peso\,GP = \frac{Gasto\,P\'ublico}{PIB} \cdot 100
\]

\begin{itemize}
	\item \textbf{Presión Fiscal:} Mide la carga tributaria efectiva sobre la economía.
\end{itemize}

\[
Presi\'on\,Fiscal = \frac{Impuestos + Cotizaciones\,Sociales}{PIB} \cdot 100
\]

\begin{itemize}
	\item \textbf{Ratio de Deuda Pública:}
\end{itemize}

\[
Ratio\,Deuda = \frac{Deuda\,P\'ublica\,Total}{PIB} \cdot 100
\]

\begin{itemize}
	\item \textbf{Ratio de Déficit/Superávit Público:}
\end{itemize}

\[
Ratio\,Saldo = \frac{Saldo\,Presupuestario}{PIB} \cdot 100
\]

